\section{Conclusion}
MitoChime offers a practical pre-assembly chimera detection pipeline for simulated mitochondrial Illumina paired-end reads using interpretable alignment-derived and sequence-based signals. The tuned Gradient Boosting classifier achieves robust performance with conservative filtering, while a CNN1D model provides higher sensitivity and stronger overall chimera detection. External assembly evaluation further shows that filtering can substantially reduce assembly graph complexity under high chimera percentage, with coverage-dependent trade-offs at low depth. In general, the Gradient Boosting model is better suited to low-coverage datasets where preserving read support is critical, while the CNN1D is more advantageous at moderate-to-high coverage where its aggressive chimera detection can be fully utilized without risking assembly fragmentation. These trade-offs may also be adjusted in the chimera detection pipeline through model selection and threshold tuning, depending on the characteristics of the input data and the specific goals of the user. Overall, these findings support MitoChime as a useful quality-control step for  improving mitochondrial genome assembly workflows, particularly in resource-constrained settings.

The main limitations of this study are its reliance on simulated reads and evaluation on a single target species. Although simulation enabled controlled benchmarking with known clean and chimeric labels, it cannot fully capture the complexity of empirical sequencing libraries, including uneven coverage, real sequencing errors, contamination, and protocol-specific artifacts. Broader validation across taxa, library preparation conditions, and real mitochondrial datasets is therefore required to establish generalizability. Future work will focus on calibration using empirical sequencing libraries, evaluation under varied sequencing error profiles, and support for additional references and model families.


